% Source-derived chapter 11: Classification of certain jacobian fibrations
% Original source: no-enriques-surface-over-the-integers
\chapter{Classification of certain jacobian fibrations}
\label{no-enriques-surface-over-the-integers-chapter-11}
\source{no-enriques-surface-over-the-integers}

\begin{remark}[Source section]
\label{no-enriques-surface-over-the-integers-chapter-11-source}
\source{no-enriques-surface-over-the-integers}
\source{no-enriques-surface-over-the-integers}
This chapter reproduces the corresponding section of the retrieved
LaTeX source, including its definitions, statements, examples, and
proofs. It has not yet been translated into Lean.
\notready
\end{remark}

% The source section title was: Classification of certain jacobian fibrations
\mylabel{Classification jacobian}

Throughout, we work over the ground field $k=\FF_2$.
Let $Y$ be a geometrically rational surface, and $\phi:Y\ra\PP^1$ be an elliptic fibration that
is relatively minimal and jacobian. It is then described by some Weierstra\ss{} equation 
$$
y^2+a_1xy+a_3y = x^3+a_2x^2+a_4x+a_6,
$$
where the coefficients $a_i\in \FF_2[t]$ are polynomials of degree $\deg(a_i)\leq i$, and the discriminant does not vanish.
The goal of this section is to classify those Weierstra\ss{} equations that might arise from
an Enriques surface with constant Picard scheme.
It turns out that there are exactly eleven such Weierstra\ss{} equations, up to change of coordinates and
automorphism of the projective line.
This is quite remarkable, because all in all there are $2^{21}=2097152$ possible Weierstra\ss{} equations.

We make the following assumptions on the surface $J$ and the jacobian fibration $\phi:J\ra\PP^1$:
\newcounter{J}\setcounter{J}{0}
\renewcommand{\theJ}{\rm \roman{J}}
\newcommand{\labelJ}[1]{\refstepcounter{J}\mylabel{#1}}
\newcommand{\refJ}[1] {{\rm (\ref{#1})}}

\begin{enumerate}
\item 
\labelJ{constant}
The proper scheme $J$ has constant   Picard scheme $\Pic_{J/k}$.
\item 
\labelJ{rational}
There is at most one rational point $a\in\PP^1$ whose fiber $J_a$   is semistable or supersingular.
\end{enumerate}

The first condition arises from Theorem \ref{jacobian constant pic}, and the second from Proposition \ref{jacobian fiber properties}.
With Propositions \ref{constant num over finite field}, \ref{canonical type rational points} and \ref{points on enriques} 
we see that   the following holds as well:

\begin{enumerate}
\setcounter{enumi}{2}
\item 
\labelJ{non-rational}
The fibers $J_\ba$ over non-rational points $a\in\PP^1$ have Kodaira $\I_0$, $\I_1$ or $\II$.
\item 
\labelJ{I0*}
There is no fiber $J_\ba$ with Kodaira symbol $\I_0^*$.
\item 
\labelJ{points}
The number of rational points is $\Card Y(\FF_2)=25$.
\end{enumerate}

For each rational point $a\in \PP^1$, let  
$r_a\geq 1$ and $n_a\geq 1$  be the respective numbers of irreducible components and rational points in $J_a$.
We start our analysis by observing that there is a large fiber, in the sense that there are many irreducible
components:

\begin{proposition}
\source{no-enriques-surface-over-the-integers}
\notready
\mylabel{six   components}
There is a rational point $a\in\PP^1$ whose fiber $J_a$ contains at least six irreducible components.
\end{proposition}

\proof
Seeking a contradiction, we   assume that $r_a\leq 5$ for the three rational points $a\in\PP^1$.
If $J_a$ is unstable we actually have $r_a\leq 3$, because the Kodaira symbol $\I_0^*$ is impossible,
and  get $n_a\leq 7$.
The same estimate holds if the fiber is smooth, because it is then isomorphic to one of the elliptic curves
$E_i$, with $1\leq i\leq 5$.
In the semistable case we get $n_a\leq 10$, and there is at most one rational point with
semistable fiber. Writing $a,b,c\in\PP^1$ for the three rational points,
we get the  contradiction $25=\Card Y(\FF_2)=n_a+n_b+n_c\leq 10 + 7 + 7= 24$.
\qed

\medskip
We shall see that this large fiber must be unstable. The next result is a  first step into this direction:

\begin{proposition}
\source{no-enriques-surface-over-the-integers}
\notready
\mylabel{unstable fiber}
There is a rational point $a\in\PP^1$ whose fiber $J_a$ is reducible and unstable.
\end{proposition}

\proof
Seeking a contradiction, we assume that all reducible fibers are semistable.
According to Proposition \ref{six components}, there is a reducible fiber with at least six components, 
say $J_a$.
It follows that $J_a$ has Kodaira symbol $\I_r$ for some $r\geq 6$,
contributing $n=2r\geq 12$ rational points. We also have the upper bound $r\leq 9$, because
of the Picard number $\rho(J)=10$.

For each of the  two other rational points $b,c\in \PP^1$, it follows that the fiber 
is either ordinary or with Kodaira symbol $\II$.
In the former case, the number of rational points is even, in the latter case it is odd.
Since $\Card Y(\FF_2)=25$ is odd, we may assume without restriction that $J_b=E_{2i}$
for $1\leq i\leq 2$, and  $J_c$ has Kodaira symbol $\II$.
This gives $25=\Card Y(\FF_2) = 2r +  2i+3$. The only possible solution is $r=9$ and $i=2$.
But by Lang's Classification \cite{Lang 2000}, the configuration of singular fibers $\II+\I_9$ does not occur.
\qed

\medskip
The generic fiber   of the jacobian fibration $\phi:J\ra\PP^1$
has a certain invariant $j(J/\PP^1)$ in the function field $ k(\PP^1)=\FF_2(t)$, which is also called
the \emph{functional $j$-invariant} and  can   be seen  as a morphism $j:\PP^1\ra\PP^1$.
It is now convenient to distinguish the cases that the invariant   is zero or non-zero.

\begin{proposition}
\source{no-enriques-surface-over-the-integers}
\notready
\mylabel{non-constant j-invariant}
Suppose the functional $j$-invariant   is non-zero.
Then there is  exactly one  unstable fiber.
If this fiber contains $r\leq 5$ irreducible components, the configuration of fibers over   rational points
must be $\III+E_4+ \I_8$. 
\end{proposition}

\proof 
First note that by Proposition \ref{unstable fiber} there is a rational point $a\in\PP^1$ whose fiber $J_a$ is unstable.
As observed by Lang, there is no other unstable fiber (\cite{Lang 2000}, beginning of Section 1).
Indeed,  a fiber $J_x$ is unstable if and only if $c_4\in \FF_2[T]$ vanishes at the point $x$ (\cite{Deligne 1975}, Proposition 5.1), and 
in our situation  $c_4=a_1^4$ with $\deg(a_1)\leq 1$.
Assume now that $J_a$   contains at most five irreducible components. Since the Kodaira symbol $\I_0^*$ does not occur,
$J_a$ has at most three irreducible components.
According to Proposition \ref{six components}, there must be a semistable fiber $J_b$ over some rational point $b\in\PP^1$,
with $6\leq m\leq 9$ irreducible components.
In light of Lang's Classification \cite{Lang 2000}, Section 2 the only possible configuration of singular fibers
is $\III+\I_6$ and $\III+\I_8$. Consequently, the fiber over the last rational point $c\in\PP^1$ is ordinary,
whence $J_c=E_{2n}$ with $1\leq n\leq 2$.
The respective number of rational points in the fiber is thus 
$n_a=5$ and $n_b=2m$ and $n_c=2n$. From $25=\Card Y(\FF_2)=n_a+n_b+n_c$ we get $m+n=10$.
The only solution is  $m=8$ and $n=2$.
\qed

\medskip
We now can state the first half our our classification result:

\begin{theorem}
\source{no-enriques-surface-over-the-integers}
\notready
\mylabel{classification non-zero j}
Suppose the functional $j$-invariant is non-zero. 
Up to coordinate changes and automorphisms of the projective line,
the fibration $\phi:J\ra\PP^1$ is given by exactly one of the following eight Weierstra\ss{} equations: 
$$
\begin{array}[c]{@{}lll}
\toprule
\text{\rm Weierstra\ss{} equation}      & \text{\rm fibers over $t=0,1,\infty$}	& \text{\rm $j$-invariant}\\
\toprule
y^2+txy+t^2y=x^3+ tx^2+t^4x+t^5(1+t)    & \I_1^*+E_4+\I_4                       & t^4\\
\midrule
y^2+txy=x^3+t^3x+t^5(1+t)               & \III^*+E_4+E_4                        & t^2/(t^2+t+1)\\
\midrule
y^2+txy=x^3+t^3x 			            & \III^*+E_4+\I_2                       & t^2\\
y^2+txy=x^3+t^2x^2+t^3x                 & \III^*+E_2+\tI_2\\
\midrule
y^2+txy=x^3+t^5                         & \II^*+E_4+\I_1                        & t\\
y^2+txy=x^3+t^2x^2+t^5                  & \II^*+E_2+\tI_1\\
\midrule
y^2+txy=x^3+tx^2+t^4x                   & \I_4^* +E_2+E_4                       & 1\\
\midrule
y^2+txy+ty=x^3+tx^2+tx                     & \III+E_4+\I_8                         & t^8\\
\bottomrule
\end{array}
$$
For the invariant $j=t^2/(t^2+t+1)$, there is an additional singular fiber, which occurs over the point $x\in\PP^1$ with $\kappa(x)\simeq\FF_4$
and has Kodaira symbol $\I_1$.
In all other cases, the fibers over the non-rational points are smooth.
\end{theorem}

\proof
Let $J_a$ be the unstable fiber, with $1\leq r_a\leq 9$ irreducible components and $n_a=2r_a+1$ rational points. 
Note that the latter is odd.
Without restriction, we may assume that $a$ is given by $t=0$. Write $b,c\in\PP^1$ for the remaining two rational points,
whence
\begin{equation}
\label{cardinality}
\source{no-enriques-surface-over-the-integers}
25=\Card Y(\FF_2) = n_a+n_b+n_c.
\end{equation}

First, assume that $r_a\geq 6$.
To start with, we consider the situation that both fibers $J_b$ and $J_c$ are smooth.
Applying an automorphism of the projective line, we may assume that the first is ordinary. Then $n_b$ is even.
By \eqref{cardinality} the number $n_c$ is also even, hence $J_c$ is also ordinary.
Write  $J_b=E_{2i}$ and $J_c=E_{2j}$. Applying another automorphism of the projective line,  we may assume
that $b$ is given by $t=1$. Recall that $r_a\geq 1$ is the number of irreducible components in $J_a$.
 From \eqref{cardinality} one infers $r_a+i+j=12$.
Using $i+j\leq 4$ we obtain $r_a\geq 8$.
Consequently, the possible configurations of   fibers over rational points are 
$$
\II^*+E_2+E_4,\quad \I_4^*+E_2+E_4,\quad \III^*+E_4+E_4,\quad \I_3^*+E_4+E_4.
$$
In the first and last cases, Lang's Classification \cite{Lang 2000} tells us that
there must be an additional singular geometric fiber, which is unique and has Kodaira symbol $\I_1$.
This uniqueness ensures that it comes from a  fiber over a rational point, contradiction.

We now analyze the case that $J_a$ has Kodaira symbol $\I_4^*$.
We are thus in Lang Case 5D from \cite{Lang 2000}, and the Weierstra\ss{} equation takes the form
$y^2+txy=x^3+td_1x^2+t^4x$, where $t\nmid d_1$.
\emph{In this context, symbols like $d_i,e_i,\ldots$ denote   polynomials in the variable $t$ of degree $\leq i$.}
Here we have the two cases $d_1=1$ and $d_1=1+t$. The former is transformed into the latter, 
by the automorphism $t\mapsto t/(1+t)$ of the projective line,  followed by the change of coordinates 
$x=u^2x'$ and $y=u^3y'$ with $u=1/(1+t)$.
For $d_1=1$ the fibers, the  Weierstra\ss{} equation and the ensuing $j$-invariant are as in the table.

Suppose now that $J_a$ has Kodaira symbol $\III^*$.
Here we are in Lang Case 7, 
and the configuration of singular geometric fibers  must
be $\III^*+\I_1+\I_1$. The Weierstra\ss{} equation is of the form
$y^2+txy+t^3c_0y=x^3+t^2d_0x^2+ t^3e_1x+t^5f_1$ with $t\nmid e_1$.
Making a change of variables we achieve $c_0=0$ and $e_1=1$.
The discriminant becomes $\Delta=t^{10}(tf_1+1)$.
Since the fiber over  the unique point $x\in\PP^1$ with residue field $\kappa(x)\simeq \FF_4$
is singular, we must have $f_1=1+t$.
The case $d_0=1$ yields $J_b=E_2$, contradiction. Thus $d_0=0$.
Again the Weierstra\ss{} equation and the ensuing $j$-invariant are as in the table.

We next consider the situation that one of the fibers over rational points is semistable.
Applying an automorphism of the projective line, we may assume that this happens over $c\in\PP^1$,
and that this point is given by $t=\infty$. Recall that $r_c\geq 1$ is the number of irreducible components
in $J_c$. The remaining fiber is ordinary,
and we write $J_b=E_{2i}$ for some $1\leq i\leq 2$, such that $n_b=2i$.
Suppose first that the semistable fiber is untwisted.
Then $n_c=2r_c$, and \eqref{cardinality} gives $r_a+i+r_c=12$. In particular, we have
$$
r_c = 12-i-r_a\geq 10-r_a.
$$
For $r_a=6$ this estimate gives $r_c\geq 4$. By Lang's Classification, the configuration of fibers over rational points
must be $\I_1^*+E_4+\I_4$. For $r_a=7$ we get  $r_a\geq 3$, and now the classification ensures that the configuration must be $\IV^*+E_4+\I_3$.
In this case, however, there is a unique additional degenerate fiber, which has Kodaira symbol $\I_1$.
The uniqueness ensures that it lies over a rational point, contradiction.
For $r_a=8$ we obtain $r_c\geq 2$, and the configuration is $\III^*+E_4+\I_2$.
Finally, for $r_a=9$ the possible configuration is $\II^*+E_4+\I_1$.
If the semistable fiber is twisted, we have  $r_c\leq 2$ and $n_c=2r_c+2$.
Now \eqref{cardinality} becomes $r_a+i+r_c=11$ and thus $r_c\geq 11-i-r_a\geq 7$, and the argument is similar.
Summing up, we have the following possibilities:
$$
\I_1^*+E_4+\I_4,\quad \III^*+E_4+\I_2,\quad\III^*+E_2+\tI_2, \quad \II^*+E_4+\I_1,\quad \II^*+E_2+\tI_1.
$$
Let me give the details for the first case: Here we are in Lang Case 5A,
with a Weierstra\ss{} equation of the form $y^2+txy+t^2c_1y=x^3+td_1x^2+t^3e_1x+t^4c_2$,
subject to the conditions  $t\nmid c_1$ and $t\nmid d_1$.
Making a change of coordinates, we may assume $c_1=1$ and $e_1=t$.
Now the discriminant becomes
$$
\Delta=t^8(t^2c_2+t^3+td_1+t^4+1+t).
$$
The second factor is a unit, hence equals one, because we assume that the semistable fiber is at $t=\infty$.
Comparing coefficients yields the solutions  $c_2=t+t^2$, $d_1=1$ and $c_2=1+t+t^2$, $d_1=1+t$.
For the latter solutions, the semistable fiber is twisted, as revealed by the Tate Algorithm at $t=\infty$.
Hence we are in the former case, and the Weierstra\ss{} equation and the resulting $j$-invariant is as in the Table.
The remaining cases are even simpler, and handled in the same fashion. We leave the details to the reader.
Summing up, this settles the case $r_a\geq 6$.

It remains to treat the case that $r_a\leq 5$. According to Proposition \ref{non-constant j-invariant}, the configuration over
the singular fibers is $\III+E_4+\I_8$, and we are in Lang Case 2A, with Weierstra\ss{} equation
$y^2+txy + tc_2y=x^3+tc_1x^2+tc_3x+t^2c_4$, with $t\nmid c_2$ and $t\nmid c_3$ and $\Delta=t^4$.
Making successive changes  of coordinates  $x=x'+R$ and $y=y'+T$ 
where the indeterminate $t$ divides $R$ and $T$, 
we may assume that $c_2=c_3=1$.
Then $1=\Delta/t^4=c_4t^4+t^3+c_1t^3+1$. It follows $c_1=1+\alpha t$ and $c_4=\alpha$
for some scalar $\alpha\in \FF_2$. In case $\alpha=1$ we have $J_b=E_2$, contradiction, hence $\alpha=0$.
\qed

\medskip
It remains to understand the situation with trivial functional $j$-invariant:

\begin{theorem}
\source{no-enriques-surface-over-the-integers}
\notready
\mylabel{classification  zero j}
Suppose the functional $j$-invariant is zero.
Up to some coordinate changes and automorphisms of the projective line,
the fibration $\phi:J\ra\PP^1$ is given by exactly one of the following three Weierstra\ss{} equations: 
$$
\begin{array}[c]{@{}ll}
\toprule
\text{\rm Weierstra\ss{} equation}		& \text{\rm fibers over $t=0,1,\infty$}\\
\toprule
y^2+t^2y=x^3+tx^2			& \I_1^*+E_5+\IV\\	
\midrule
y^2+t^2y=x^3+t^3x 			& \IV^*+E_5+\III\\
\midrule
y^2+t^2y=x^3 			& \IV^*+E_3+\IV\\
\bottomrule
\end{array}
$$
In all three cases, the fibers over non-rational points $x\in\PP^1$ are smooth.
\end{theorem}

\proof
There are neither ordinary nor semistable fibers, because the  map $j:\PP^1\ra\PP^1$ misses
the values $t=1$ and $t=\infty$. In particular, each smooth fiber  over a rational point 
is isomorphic to one of the elliptic curves $E_1$, $E_3$ or $E_5$. So there are
at least two rational points $a\neq c$ on $\PP^1$ such that $J_a$ and $J_c$ are unstable.
By Lang's Classification, there are at most three singular geometric fibers. If present, 
the third singular geometric fiber must be Galois invariant, whence occurs over
the remaining rational point $b\in\PP^1$. Thus all fibers over non-rational points are smooth.

Suppose first that all fibers of $\phi:J\ra\PP^1$ are reduced, such that the Kodaira symbols for
the singular fibers are $\II$, $\III$ or $\IV$. Then each fiber over  a rational point contains
at most seven rational points, which gives the contradiction $25=\Card Y(\FF_2)\leq 3\cdot 7$.
Thus there is a non-reduced rational fiber.

Recall that $J_a$ and $J_c$ are unstable. Let $r_a,r_c\geq 1$ be the respective number of irreducible components.
After applying an automorphism of the projective line, we may assume that $J_a $  is non-reduced. 
By Lang's Classification (\cite{Lang 2000}, page 5829), the remaining fiber $J_b$ must be smooth 
(one referee pointed out that this also follows from the fact that the valuation of the discriminant
in characteristic $p=2$ is at least four at an additive fiber, and at least eight at a non-reduced fiber).
Write $J_b=E_{2i-1}$ for some $1\leq i\leq 3$.
Then $25=\Card Y(\FF_2) = (2r_a+1) +(2i-1) + (2r_c+1) $, or equivalently $r_a+i+r_c=12$.
Since $J_a$ is non-reduced and the Kodaira symbol $\I_0^*$ does not occur, we have $6\leq r_a\leq 9$.
In particular, $r_c\leq 5$, so the unstable fiber $J_c$ is reduced, such that actually $1\leq r_c\leq 3$.
The case $r_a=6$ implies $r_c=i=3$. In case $r_a=7$ we get $r_c=3$, $i=2$ or $r_c=2$, $i=3$.
Going through Lang's Classification, we see that these are the only possibilities,
which already give the second column of the table.

Suppose the configuration of fibers over rational points is $\I_1^*+E_5+\IV$.
Then we are in Lang Case 13A, and the Weierstra\ss{} equation takes the form
$y^2+t^2c_1y=x^3+td_1x^2+t^3e_1x+t^4d_2$ with $t\nmid d_1$ and $t\nmid c_1$.
The discriminant is $\Delta=t^8c_1^4$, thus $c_1=1$. Making a change of coordinate $y'=y+sx$,
we achieve $d_1=1$. Making a further change of coordinates, we may assume
that $e_1$ has no linear term, and that $d_2$ has neither constant nor quadratic term.
Now the Weierstra\ss{} equation becomes $y^2+t^2y=x^3+tx^2+t^3e_0x+t^5d_0$.
One checks that only $e_0=d_0=0$ yields $J_b=E_5$.

Next, suppose the configuration is $\IV^*+E_5+ \III$.
Now we are in Lang Case 14, with $y^2+t^2c_1y=x^3+t^2d_0x^2+t^3e_1x+t^4d_2$ subject to
the condition $t\nmid c_1$. The discriminant is $\Delta=t^8c_1^4$, whence $c_1=1$.
Making a suitable change of variables $y=y'+sx'$ we   achieve $d_0=0$.
As in the preceding paragraph, we put the Weierstra\ss{} equation into
the form $y^2+t^2y=x^3+t^3e_0x+t^5d'_0$.
Only the case $e_0=1$ and $d'_0=0$ yields $J_b=E_5$.

Finally, consider the configuration $\IV^*+E_3+\IV$.
We argue as in the preceding paragraph to reduce to the very same   Weierstra\ss{} equation.
Now only $e_0=d'_0=0$ yields the fiber $J_b=E_3$. 
\qed


%===========================================================
